\newcommand{\tipoRequerimento}{Solicitação de Relatório de Dívidas}

\section{FUNDAMENTAÇÃO LEGAL DA REQUISIÇÃO}

\begin{enumerate}[resume]
\item
  Para embasar a requisição, utiliza-se as seguintes disposições da Resolução Normativa nº 1.000/2021 (REN nº 1.000/2021) da ANEEL:
  \begin{enumerate}
  \item
    Demonstrativos e Memória de Cálculo (Art. 472, § 1º): A distribuidora tem a obrigação de fornecer ao município os demonstrativos detalhados e a memória de cálculo do faturamento da iluminação pública. Prazos de Atendimento (Art. 408):
    \begin{enumerate}
      \item  
      5 dias úteis para respostas que não exigem visita técnica.
      \item 
      10 dias úteis para os demais casos.
    \end{enumerate}
  \item 
  Prazo Geral (Art. 409): Caso não haja um prazo específico na regulamentação da ANEEL para uma solicitação, a distribuidora deve atendê-la em até 30 dias.
  \end{enumerate}
\end{enumerate}

\section{DOS REQUERIMENTOS}

\begin{enumerate}[resume]
\item Solicita-se o envio do \textbf{Relatório de Dívidas} do município, contendo no mínimo as seguintes informações.
\begin{enumerate}
  \item Número da unidade consumidora;
  \item Origem do débito (consumo, TOI censo ou TOI comum);
  \item Data do faturamento;
  \item Data do vencimento;
  \item Valor do débito;
  \item Número da fatura;
  \item Se a fatura consta judicializada;
\end{enumerate}
\item Se possível, apresentar as informações solicitadas acima por meio de planilha em \textbf{Excel}.
\item Atender a requisição no prazo máximo de 5 (cinco) dias úteis, conforme estabelece art. 408, inciso I, da REN 1.000/2021
\item Além disso, informar se há negociações em andamento entre a prefeitura e a concessionária sobre acordos ou contestações de valores nas contas de luz.
\end{enumerate}

